%% kapa: exam=1 round=33 year=2022 period=1 form=A subject=real_estate subject_ko=부동산학원론 seq=13 number=93
%% source: Q-Net 감정평가사 기출문제 (cst00309) artlSeq=5209997 file=raw/1008/5209997/2022년도 제33회 감정평가사 1차시험 1교시(민법,경제학,부동산학).pdf
%% conversion: kordoc 4.24.51 -> tools/gen1.py (auto); review=proofread-sonnet; flags=none
\begin{question}{93}
A투자안의 현금흐름이다. 추가투자가 없었을 때의 NPV(ㄱ)와 추가투자로 인한 NPV증감(ㄴ)은? (단, 0기 기준이며, 주어진 자료에 한함)
\par{\small\centering\begin{tabular}{|l|c|c|c|c|}\hline
구분 & 0기 & 1기 & 2기 & 3기 \tabularnewline \hline
초기투자 & (1억원) & & & \tabularnewline \hline
NOI & & 4천만원 & 3천만원 & 4천만원 \tabularnewline \hline
추가투자 & & & (5천만원) & \tabularnewline \hline
추가투자에 따른 NOI증감 & & & +3천만원 & +4천만원 \tabularnewline \hline
현가계수 & & 0.952 & 0.906 & 0.862 \tabularnewline \hline
\end{tabular}\par}
\begin{choices}
  \item ㄱ: $-$260,000원, ㄴ: +16,360,000원
  \item ㄱ: $-$260,000원, ㄴ: +17,240,000원
  \item ㄱ: $-$260,000원, ㄴ: +18,120,000원
  \item ㄱ: +260,000원, ㄴ: +16,360,000원
  \item ㄱ: +260,000원, ㄴ: +17,240,000원
\end{choices}
\end{question}
\answer{1}
